\section{引言：一个被修辞掩盖的问题}\label{sec:intro}

\subsection{问题的提法}

交易策略的通俗分类把趋势跟踪与均值回归摆在对立的两端：前者顺势加仓、追涨杀跌，信奉“趋势会延续”；后者逆势布局、买跌卖涨，信奉“偏离终将修复”。坊间叙述往往把这种对立上升为世界观之争——一种相信市场会走远，另一种相信市场会回头。

这类叙述有一个不易察觉的问题：它把\emph{交易动作}的差异误当成了\emph{信息来源}的差异。追涨与买跌确实是相反的动手方式，但两者据以决策的原料完全相同——都是资产过去一段时间的价格变化。趋势策略把这段历史当作延续的证据，反转策略把同一段历史当作过度的证据。二者从未使用不同的信息集，只是对同一份信息给出了相反的读数。

由此可以给出本文的核心论断：

\begin{quote}
\itshape
趋势跟踪与均值回归不是两套独立的世界观，而是同一个预测信号的两种符号选择。
把它们区分开的，既不是信息集、也不是数学工具，而是关于\emph{收益自相关符号}的一个假设。
\end{quote}

这个论断如果成立，会带来一连串可检验的推论：两条文献脉络应当在历史上共享源头；同一份数据的自相关函数应当在不同时间尺度上同时呈现正负两种符号；两个策略的收益分布应当在偏度上呈现系统性的镜像；而最直接的一条是——用同一资产池、同一成本模型、仅仅翻转信号符号，应当能复现出两者的绩效分野。

本文的任务就是把这个论断拆开、论证、并放到数据上检验。

\subsection{为什么这件事值得认真做}

否认“两策略本质对立”的说法，常被误解为一种和解式的修辞。恰恰相反：把它讲清楚是有实际后果的。

其一，\textbf{它改变了参数选择的方式}。如果两策略共享同一信号，那么决定用哪一个的关键变量就不是“你相信哪种市场观”，而是\emph{你打算在哪个时间尺度上持仓}。Bachelier 以来的收益序列自相关研究反复指向同一个经验规律：短周期与长周期呈现负自相关，中周期呈现正自相关。这意味着趋势与反转的归属由持有期决定，而不是由信念决定。把选策略的问题转化为选时间尺度的问题，是一个更接近可操作层面的提法。

其二，\textbf{它改写了风险预算的对象}。两者的收益分布并非对称的镜像。趋势策略的收益右偏、伴随罕见的深度左尾（动量崩溃），本质上是持有一个隐含的看涨期权结构；反转策略的收益左偏、依赖价格收敛，本质上是承担了卖方期权式的尾部风险\cite{funghsieh2001,danielmoskowitz2016}。如果只比较夏普比率而忽略偏度，就会把一个“平时赚钱、偶尔跳崖”的策略与一个“常态波动、极端时刻被资助”的策略当成同类。执行层面的止损、杠杆、再加仓规则，必须建立在对这种偏度差异的认识之上。

其三，\textbf{它暴露了共同的脆弱点}。两策略的发现在方法论上面对同一类质疑：多重检验与数据窥探。Sullivan 等(1999)对 7846 条技术规则的检验、White(2000)的 Reality Check、Harvey 等(2016)提出的更高 $t$ 值门槛，指向的是同一个问题——在一个被反复搜索的数据集上，“显著”是廉价的\cite{sullivan1999,white2000,hlz2016}。趋势与反转都难逃这一审查；把二者放入同一审计框架，比分别辩护更有说服力。

\subsection{比较的判据}

要比较两个策略，先要定义“比什么”。若不加限定地比较，结论会轻易退化为不可比的数字竞赛：一个是月频、低换手、跨资产的组合，另一个是日频、高换手、单资产的价差交易，说“前者年化更高”没有意义。

本文确立六个维度的比较判据，后文各章围绕它们展开：

\begin{enumerate}[leftmargin=2.2em]
  \item \textbf{信息与数学结构}：用什么变量预测什么量？预测函数的形状与符号如何？这是“本质”层面。
  \item \textbf{统计对象}：策略押注的究竟是哪一个可检验量——自相关系数、方差比、协整关系、还是半衰期？这决定了策略是否“先有入场券”。
  \item \textbf{收益分布}：偏度、峰度、以及极端市场中的表现（凸性 vs 凹性）。
  \item \textbf{执行代价}：换手量级、对交易成本的敏感度、容量与拥挤。
  \item \textbf{失效模式}：什么情况下会亏钱？亏钱是渐进的还是突发的？失败是可止损的还是“越止损越糟”的？
  \item \textbf{证据质量}：支持各自的核心证据在数据窥探、样本外衰减、发表后衰减上表现如何。
\end{enumerate}

\subsection{研究方法与证据纪律}

本文属于证据型综述，方法上遵循三条纪律。

\textbf{一是全部引用可回溯至原始来源}。文献的题名、作者、年份、期刊与 DOI 均通过 OpenAlex、Crossref、arXiv 官方接口或出版社页面逐一核对，不以记忆填写。凡接口返回与记忆冲突者，以接口为准。少数无法通过公开接口确证的经典（如 Alexander 1961 的早期工作），明确标注为待复核，并在正文中避免依赖其具体数值。

\textbf{二是区分“机制解释”与“实证事实”}。行为学解释（反应不足、过度反应）与风险补偿解释（时变风险溢价）在文献中并存且各有支持者，本文不做仲裁，而是并列呈现，并在结论置信度一章逐条给出分级\cite{cochrane2011,bsv1998,dhs1998,hongstein1999}。

\textbf{三是自备一组统一口径的对照实证}。文献中的绩效数字来自不同样本、不同成本假设、不同资产池，横向比较的可靠性有限。为检验“符号翻转”这一核心预测，本文在中国期货主力连续合约与 A 股两个资产池上重建了一套对照回测：两策略共享同一资产池、同一调仓频率、同一成本模型与同一评估指标，唯一区别是信号符号（详见第~\ref{sec:empirical}~节）。这组实证不宣称发现新的alpha，它的作用是\emph{把文献论断放到一个可控的变量对照下}。

\subsection{边界与结构}

本文的比较限于\emph{可交易的收益自相关策略}。宏观计量意义上的均值回归（如购买力平价的长期收敛检验）与非交易性的价格发现研究，仅在涉及机制解释时引用，不作为策略比较的对象。此外，本文讨论的市场证据以发达市场文献为主，中国市场部分主要依赖本文自建的对照实证与国内可得数据，样本长度与制度差异构成外推边界，将在第~\ref{sec:audit}~节与结论中明确。

全文结构如下。第~\ref{sec:genealogy}~节并置两条脉络的编年史；第~\ref{sec:unity}~节给出统一框架；第~\ref{sec:duality}~节论证七组本质差异；第~\ref{sec:matrix}~节用统一比较矩阵汇总证据；第~\ref{sec:mechanism}~节对照机制解释；第~\ref{sec:execution}~节落到执行差异；第~\ref{sec:empirical}~节是本文自建的对照实证；第~\ref{sec:audit}~节做污染与不可比审计；此后依次是结论置信度分级、可复现实验路线图、以及研究现状与遗留课题。
